\chapter{Kardinalität und Potenzmenge}

\section{Die beweisbare Vorwärtsrichtung}

\FormulaThmDeltaK[Transport einer Bijektion auf die Potenzmengen]
{A\approx B\ \vdash\ \mathcal P(A)\approx\mathcal P(B)}
{B48PowerTransport}{
  \DeltaRow{Mengen}{A,B}
  \DeltaRow{Bijektion}{f:A\longrightarrow B}
}
\begin{tabproof}
 \proofstep{1}{f:A\longrightarrow B\text{ bijektiv}}{\rA}
 \proofstep{1}{\begin{aligned}
 F:\mathcal P(A)&\longrightarrow\mathcal P(B),
 &F(X)&=f[X],\\
 H:\mathcal P(B)&\longrightarrow\mathcal P(A),
 &H(Y)&=f^{-1}[Y]
 \end{aligned}}
 {\BXLVIIIWhy{Bild und Urbild sind Teilmengen der richtigen Grundmengen}}
 \proofstep{1}{x\in H(F(X))\ \Longleftrightarrow\
 f(x)\in f[X]\ \Longleftrightarrow\ x\in X}
 {\BXLVIIIWhy{Injektivität von \(f\)}}
 \proofstep{1}{F(H(Y))=Y}
 {\BXLVIIIWhy{Jedes \(y\in Y\) hat wegen Surjektivität ein Urbild;
 die andere Inklusion folgt aus der Urbilddefinition}}
 \proofstep{1}{H\circ F=\operatorname{id}_{\mathcal P(A)},\quad
 F\circ H=\operatorname{id}_{\mathcal P(B)}}
 {\BXLVIIIWhy{Extensionalität}}
 \proofstep{1}{\mathcal P(A)\approx\mathcal P(B)}
 {\BXLVIIIWhy{\(F\) und \(H\) sind inverse Bijektionen}}
\end{tabproof}
\noindent
Dieser Beweis benötigt kein Auswahlaxiom und gilt auch für leere
Mengen. Er erklärt außerdem genau, was das Wort
\enquote{induziert} bei der Potenzmengenabbildung bedeutet.

\section{Cantors Satz}

\FormulaThmDeltaK[Eine Menge ist kleiner als ihre Potenzmenge]
{\forall A\quad |A|<|\mathcal P(A)|}
{B48Cantor}{
  \DeltaRow{Menge}{A}
}
\begin{tabproof}
 \proofstep{}{\iota:A\longrightarrow\mathcal P(A),\quad
 \iota(a)=\{a\}\text{ ist injektiv}}
 {\BXLVIIIWhy{Gleichheit von Singletons}}
 \proofstep{2}{f:A\longrightarrow\mathcal P(A)\text{ surjektiv}}{\rA}
 \proofstep{2}{D=\{a\in A:a\notin f(a)\}\in\mathcal P(A)}
 {\BXLVIIIWhy{Aussonderung}}
 \proofstep{2}{\exists d\in A\ (f(d)=D)}
 {\BXLVIIIWhy{Surjektivität}}
 \proofstep{2}{d\in D\ \Longleftrightarrow\
 d\notin f(d)\ \Longleftrightarrow\ d\notin D}
 {\BXLVIIIWhy{Definition von \(D\), Ersetzung von \(f(d)\)}}
 \proofstep{}{\BXLVIIIText{Es gibt keine Surjektion von \(A\)
 auf \(\mathcal P(A)\), insbesondere keine Bijektion.
 Mit der Singleton-Injektion ergibt dies die strikte
 Mächtigkeitsungleichung.}}
 {\BXLVIIIWhy{Widerspruch, 2 entladen}}
\end{tabproof}
\noindent
Die Diagonalargumentation selbst gilt bereits ohne Auswahl.
In der Kardinalnotation von ZFC folgt insbesondere
\(\kappa<2^\kappa\) für jede Kardinalzahl \(\kappa\).

\section{Wann die Rückrichtung gilt}

\FormulaThmDeltaK[Potenzmengenkürzung im endlichen Fall]
{\begin{gathered}
 A\text{ endlich},\quad\mathcal P(A)\approx\mathcal P(B)\\
 \vdash\ B\text{ endlich und }A\approx B
 \end{gathered}}
{B48FiniteCancellation}{
  \DeltaRow{Mengen}{A,B}
}
\begin{tabproof}
 \proofstep{1}{|A|=n<\omega,\quad
 \mathcal P(A)\approx\mathcal P(B)}{\rA}
 \proofstep{1}{|\mathcal P(A)|=2^n}
 {\BXLVIIIWhy{Charakteristische Funktionen;
 je Element zwei Möglichkeiten}}
 \proofstep{1}{\BXLVIIIText{Die Singleton-Injektion von \(B\)
 in \(\mathcal P(B)\) zeigt, dass auch \(B\) endlich ist.
 Setze \(|B|=m\). Dann gilt \(2^n=2^m\).}}
 {\BXLVIIIWhy{Teilmengen endlicher Mengen sind endlich}}
 \proofstep{1}{n=m}
 {\BXLVIIIWhy{\(2^{k+1}=2\cdot2^k>2^k\);
 strenge Monotonie auf den natürlichen Zahlen}}
 \proofstep{1}{A\approx B}{\BXLVIIIWhy{Gleiche endliche Mächtigkeit}}
\end{tabproof}

\FormulaThmDeltaK[Kürzung als striktes Wachstum der Potenzfunktion]
{\BXLVIIIPC\ \Longleftrightarrow\
 \forall\kappa<\lambda\text{ Kardinalzahlen}\ (2^\kappa<2^\lambda)}
{B48PCStrictGrowth}{
  \DeltaRow{Theorie}{\BXLVIIIZFC}
}
\begin{tabproof}
 \proofstep{1}{\BXLVIIIPC,\quad\kappa<\lambda}{\rA}
 \proofstep{1}{2^\kappa\le2^\lambda}
 {\FormulaRefAuto{B48ExponentLaws}[thm]}
 \proofstep{1}{\BXLVIIIText{Gleichheit würde eine Bijektion
 zwischen den Potenzmengen von \(\kappa,\lambda\) und
 nach PK eine Bijektion zwischen \(\kappa,\lambda\)
 liefern. Das widerspricht ihrer Verschiedenheit als
 Kardinalzahlen. Also gilt strikte Ungleichheit.}}
 {\BXLVIIIWhy{Definition von PK und Kardinalzahlen}}
 \proofstep{4}{\BXLVIIIText{Die Potenzfunktion wachse strikt
 auf den Kardinalzahlen, und
 \(\mathcal P(A)\approx\mathcal P(B)\).}}{\rA}
 \proofstep{4}{2^{|A|}=2^{|B|}}
 {\FormulaRefAuto{B48ExponentLaws}[thm]}
 \proofstep{4}{\BXLVIIIText{Wäre \(|A|\ne|B|\), wäre nach dem
 Kardinalvergleich eine der beiden Kardinalzahlen kleiner.
 Striktes Wachstum widerspräche der vorigen Gleichheit.
 Also \(|A|=|B|\) und \(A\approx B\).}}
 {\BXLVIIIWhy{Ordinalvergleich für Kardinale}}
 \proofstep{}{\BXLVIIIPC\ \Longleftrightarrow\
 \forall\kappa<\lambda\ (2^\kappa<2^\lambda)}
 {\BXLVIIIWhy{Beide Richtungen, Voraussetzungen entladen}}
\end{tabproof}

\FormulaThmDeltaK[GCH impliziert die globale Potenzmengenkürzung]
{\BXLVIIIZFC\vdash\BXLVIIIGCH\longrightarrow\BXLVIIIPC}
{B48GCHImpliesPC}{
  \DeltaRow{Aussage}{\BXLVIIIPC\text{ aus Kapitel 1}}
}
\begin{tabproof}
 \proofstep{1}{\BXLVIIIGCH,\quad\kappa<\lambda}{\rA}
 \proofstep{1}{\BXLVIIIText{Sind beide Kardinale endlich,
 so gilt \(2^\kappa<2^\lambda\) nach der endlichen
 Arithmetik. Ist nur \(\kappa\) endlich, so ist
 \(2^\kappa\) endlich, \(2^\lambda\) dagegen unendlich.}}
 {\BXLVIIIWhy{Endliche Potenzen und Singleton-Injektion}}
 \proofstep{1}{\begin{aligned}
 \kappa\ge\aleph_0\quad&\Longrightarrow\\
 2^\kappa=\kappa^+&\le\lambda<\lambda^+=2^\lambda
 \end{aligned}}
 {\BXLVIIIWhy{GCH und Definition der Nachfolgerkardinale}}
 \proofstep{1}{\forall\kappa<\lambda\ (2^\kappa<2^\lambda)}
 {\BXLVIIIWhy{Vollständige Fallunterscheidung}}
 \proofstep{1}{\BXLVIIIPC}
 {\FormulaRefAuto{B48PCStrictGrowth}[thm]}
\end{tabproof}

\section{Ein Modell mit gleicher Potenzmenge bei verschiedenen Kardinalen}

\FormulaThmDeltaK[Das benötigte Plateau der Potenzfunktion]
{\begin{gathered}
 M\models\BXLVIIIZFC+\BXLVIIIGCH,\quad |M|\le\aleph_0,\\
 \mathbb P=\operatorname{Add}(\omega,\aleph_2)^M,\quad
 G\text{ \(M\)-generisch}\\
 \vdash\ M[G]\models
 \BXLVIIIZFC+\bigl(2^{\aleph_0}=2^{\aleph_1}=\aleph_2\bigr)
 \end{gathered}}
{B48PlateauModel}{
  \DeltaRow{Konstruktion}{\lambda=(\aleph_2)^M}
}
\begin{tabproof}
 \proofstep{1}{M\models\BXLVIIIZFC+\BXLVIIIGCH,\quad
 \mathbb P=\operatorname{Add}(\omega,\lambda)^M}{\rA}
 \proofstep{1}{M[G]\models\BXLVIIIZFC}
 {\FormulaRefAuto{B48ForcingZFC}[thm]}
 \proofstep{1}{\BXLVIIIText{In \(M\) hat \(\mathbb P\)
 Mächtigkeit \(\lambda\) und erfüllt ccc.
 Daher sind \(\aleph_0,\aleph_1,\aleph_2\) in der
 Erweiterung dieselben Kardinale unter der kanonischen
 Einbettung.}}
 {\BXLVIIIWhy{Cohen-ccc und Kardinalerhaltung}}
 \proofstep{1}{\begin{aligned}
 M\models\lambda^{\aleph_1}
 &=(2^{\aleph_1})^{\aleph_1}\\
 &=2^{\aleph_1\cdot\aleph_1}
 =2^{\aleph_1}=\lambda
 \end{aligned}}
 {\BXLVIIIWhy{GCH im Ausgangsmodell; Exponentengesetze und Produkte}}
 \proofstep{1}{M\models(\lambda^{\aleph_0})^{\aleph_1}
 =\lambda^{\aleph_0\cdot\aleph_1}=\lambda}
 {\BXLVIIIWhy{Exponentengesetz und vorige Zeile}}
 \proofstep{1}{M[G]\models 2^{\aleph_1}\le\lambda}
 {\FormulaRefAuto{B48NiceNameBound}[thm]}
 \proofstep{1}{M[G]\models\lambda\le2^{\aleph_0}}
 {\FormulaRefAuto{B48CohenReals}[thm]}
 \proofstep{1}{M[G]\models
 \lambda\le2^{\aleph_0}\le2^{\aleph_1}\le\lambda}
 {\BXLVIIIWhy{Monotonie der Potenzfunktion}}
 \proofstep{1}{M[G]\models
 2^{\aleph_0}=2^{\aleph_1}=\aleph_2}
 {\BXLVIIIWhy{Antisymmetrie des Kardinalvergleichs}}
\end{tabproof}

\noindent
Nur viele neue reelle Zahlen hinzuzufügen würde noch nicht genügen:
Das gäbe lediglich eine untere Schranke für \(2^{\aleph_0}\).
Erst die Zählung der schönen Namen liefert die obere Schranke
für \(2^{\aleph_1}\) und damit die tatsächlich benötigte Gleichheit.
Auch die Erhaltung von \(\aleph_1,\aleph_2\) ist ein eigener,
zuvor bewiesener Bestandteil des Arguments.

\FormulaThmDeltaK[Das Plateau widerlegt die Potenzmengenkürzung]
{\BXLVIIIZFC\vdash
 (2^{\aleph_0}=2^{\aleph_1})\longrightarrow\neg\BXLVIIIPC}
{B48PlateauRefutesPC}{
  \DeltaRow{Zeugen}{A=\omega,\quad B=\omega_1}
}
\begin{tabproof}
 \proofstep{1}{2^{\aleph_0}=2^{\aleph_1}}{\rA}
 \proofstep{1}{|\mathcal P(\omega)|=|\mathcal P(\omega_1)|}
 {\FormulaRefAuto{B48ExponentLaws}[thm]}
 \proofstep{1}{\mathcal P(\omega)\approx\mathcal P(\omega_1)}
 {\BXLVIIIWhy{Bedeutung gleicher Kardinalität}}
 \proofstep{}{|\omega|=\aleph_0<\aleph_1=|\omega_1|}
 {\FormulaRefAuto{B48Hartogs}[thm]}
 \proofstep{1}{\begin{gathered}
 \mathcal P(\omega)\approx\mathcal P(\omega_1)\\
 \land\quad\omega\not\approx\omega_1
 \end{gathered}}
 {\BXLVIIIWhy{Vorige Zeilen}}
 \proofstep{1}{\neg\BXLVIIIPC}
 {\BXLVIIIWhy{Explizites Gegenbeispiel zur universellen Aussage}}
\end{tabproof}

\section{Unabhängigkeit der Potenzmengenkürzung}

\FormulaThmDeltaK[Beide Seiten der Potenzmengenkürzung sind relativ konsistent]
{\begin{gathered}
 \BXLVIIICon(\BXLVIIIZFC)\ \Longrightarrow\\
 \BXLVIIICon(\BXLVIIIZFC+\BXLVIIIPC)
 \ \land\
 \BXLVIIICon(\BXLVIIIZFC+\neg\BXLVIIIPC)
 \end{gathered}}
{B48PCBothConsistent}{
  \DeltaRow{Voraussetzung}{\BXLVIIICon(\BXLVIIIZFC)}
}
\begin{tabproof}
 \proofstep{1}{\BXLVIIICon(\BXLVIIIZFC)}{\rA}
 \proofstep{1}{\BXLVIIICon(\BXLVIIIZFC+\BXLVIIIGCH)}
 {\FormulaRefAuto{B48GodelConsistency}[thm]}
 \proofstep{1}{M\models\BXLVIIIZFC+\BXLVIIIGCH,\quad |M|\le\aleph_0}
 {\FormulaRefAuto{B48Completeness}[thm]}
 \proofstep{1}{M\models\BXLVIIIPC}
 {\FormulaRefAuto{B48GCHImpliesPC}[thm]}
 \proofstep{1}{\BXLVIIICon(\BXLVIIIZFC+\BXLVIIIPC)}
 {\FormulaRefAuto{B48ModelConsistency}[thm]}
 \proofstep{1}{G\text{ ist \(M\)-generisch für }
                       \operatorname{Add}(\omega,\aleph_2)^M}
 {\FormulaRefAuto{B48GenericExists}[thm]}
 \proofstep{1}{M[G]\models\BXLVIIIZFC+
                 (2^{\aleph_0}=2^{\aleph_1}=\aleph_2)}
 {\FormulaRefAuto{B48PlateauModel}[thm]}
 \proofstep{1}{M[G]\models\neg\BXLVIIIPC}
 {\FormulaRefAuto{B48PlateauRefutesPC}[thm]}
 \proofstep{1}{\BXLVIIICon(\BXLVIIIZFC+\neg\BXLVIIIPC)}
 {\FormulaRefAuto{B48ModelConsistency}[thm]}
 \proofstep{1}{\BXLVIIICon(\BXLVIIIZFC+\BXLVIIIPC)
 \land\BXLVIIICon(\BXLVIIIZFC+\neg\BXLVIIIPC)}
 {\BXLVIIIWhy{Konjunktionseinführung}}
\end{tabproof}

\FormulaThmDeltaK[Unabhängigkeit der globalen Potenzmengenkürzung]
{\begin{gathered}
 \BXLVIIICon(\BXLVIIIZFC)\ \Longrightarrow\\
 (\BXLVIIIZFC\nvdash\BXLVIIIPC)
 \ \land\ (\BXLVIIIZFC\nvdash\neg\BXLVIIIPC)
 \end{gathered}}
{B48PCIndependence}{
  \DeltaRow{Aussage}{
  \BXLVIIIPC:\ \forall A,B\
       (\mathcal P(A)\approx\mathcal P(B)\to A\approx B)}
}
\begin{tabproof}
\proofstep{1}{\BXLVIIICon(\BXLVIIIZFC)}{\rA}
 \proofstep{1}{\begin{gathered}
   \BXLVIIICon(\BXLVIIIZFC+\BXLVIIIPC)\\
   \land\ \BXLVIIICon(\BXLVIIIZFC+\neg\BXLVIIIPC)
   \end{gathered}}
 {\FormulaRefAuto{B48PCBothConsistent}{1}}
 \proofstep{1}{\begin{gathered}
   (\BXLVIIIZFC\nvdash\BXLVIIIPC)\\
   \land\ (\BXLVIIIZFC\nvdash\neg\BXLVIIIPC)
   \end{gathered}}
 {\FormulaRefAuto{B48IndependenceCriterion}{2}}
 \proofstep{}{\begin{gathered}
   \BXLVIIICon(\BXLVIIIZFC)\longrightarrow\\
   ((\BXLVIIIZFC\nvdash\BXLVIIIPC)\land
    (\BXLVIIIZFC\nvdash\neg\BXLVIIIPC))
   \end{gathered}}
 {\RuleRef{RI}{\rightarrow I}\,1\!:\!3}
\end{tabproof}

\noindent
Damit ist die in diesem Kapitel gestellte Rückrichtungsfrage
entschieden: Die universelle Rückrichtung ist, unter der
Konsistenzvoraussetzung, weder ein Satz noch widerlegbar in ZFC.
Das berührt weder den endlichen Fall noch den Satz über eine
vorgegebene induzierte Potenzmengenabbildung.
